% Physical database tables, version B (29 tables, redrawn ERD). Generated by ch6t2/gen/gen2.py.
% Type, Length, Not Null and keys: Backend_API/Migrations/AppDbContextModelSnapshot.cs;
% Default: C# property initialisers in Backend_API/Models (the model has no HasDefaultValue).
SmartExam keeps its data in 29 PostgreSQL tables, one for each entity of the ERD (Figures~\ref{fig:erd_diagram} and~\ref{fig:erd_diagram_b}). For each table, one table gives its fields with type, length, default value, whether they may be empty and its keys.

\edef\BTsavedES{\the\emergencystretch}\emergencystretch=2em
\par\ifdim\dimexpr\pagegoal-\pagetotal\relax<230pt\newpage\fi
\subsection{User Table}
Table~\ref{tab:db_user} gives the fields of the Users table, which stores every account in the system for all five roles. Most other tables point to it, and the email and roll number of an account are unique.

{\small
\begin{longtable}{>{\raggedright\arraybackslash}p{2.8cm}>{\raggedright\arraybackslash}p{3.7cm}>{\raggedright\arraybackslash}p{1.95cm}>{\centering\arraybackslash}p{1.4cm}>{\raggedright\arraybackslash}p{1.95cm}>{\centering\arraybackslash}p{1.5cm}}
\caption{User Table}\label{tab:db_user}\\
\hline\hline
\textbf{Field Name} & \textbf{Description} & \textbf{Type} & \textbf{Length} & \textbf{Default} & \textbf{Not Null} \\ \hline
\endfirsthead
\multicolumn{6}{c}{\tablename~\thetable{} -- continued from previous page}\\
\hline\hline
\textbf{Field Name} & \textbf{Description} & \textbf{Type} & \textbf{Length} & \textbf{Default} & \textbf{Not Null} \\ \hline
\endhead
\hline \multicolumn{6}{r}{Continued on next page}\\
\endfoot
\hline\hline
\endlastfoot
UserId & Unique id of the user & uuid & -- & Primary Key & Yes \\*
Name & Display name & Text & -- & -- & Yes \\*
Email & Email address, used to log in (unique) & Text & -- & -- & Yes \\
RollNumber & University roll number, stored in upper case; students may log in with it (unique when set) & Varchar & 32 & -- & No \\
PasswordHash & Salted PBKDF2 hash of the password & Text & -- & -- & Yes \\
Salt & Random salt used for the hash & Text & -- & -- & Yes \\*
Role & SuperAdmin, Admin, Teacher, Invigilator or Student & Text & -- & -- & Yes \\*
IsActive & False when the account is deactivated & Boolean & -- & true & Yes \\*
CreatedAt & When the account was created & Timestamp & -- & Current time & Yes \\
\end{longtable}
}

\par\ifdim\dimexpr\pagegoal-\pagetotal\relax<230pt\newpage\fi
\subsection{User Log Table}
Table~\ref{tab:db_userlog} gives the fields of the AuditLogs table, the log of security-relevant and staff actions. Each row holds the hash of the row before it, so a changed or deleted row can be detected.

{\small
\begin{longtable}{>{\raggedright\arraybackslash}p{2.8cm}>{\raggedright\arraybackslash}p{3.7cm}>{\raggedright\arraybackslash}p{1.95cm}>{\centering\arraybackslash}p{1.4cm}>{\raggedright\arraybackslash}p{1.95cm}>{\centering\arraybackslash}p{1.5cm}}
\caption{User Log Table}\label{tab:db_userlog}\\
\hline\hline
\textbf{Field Name} & \textbf{Description} & \textbf{Type} & \textbf{Length} & \textbf{Default} & \textbf{Not Null} \\ \hline
\endfirsthead
\multicolumn{6}{c}{\tablename~\thetable{} -- continued from previous page}\\
\hline\hline
\textbf{Field Name} & \textbf{Description} & \textbf{Type} & \textbf{Length} & \textbf{Default} & \textbf{Not Null} \\ \hline
\endhead
\hline \multicolumn{6}{r}{Continued on next page}\\
\endfoot
\hline\hline
\endlastfoot
LogId & Unique id of the record & uuid & -- & Primary Key & Yes \\*
ActorId & User who acted; null for system events (plain id, no foreign key) & uuid & -- & -- & No \\*
EventType & Action, e.g. LOGIN\_SUCCESS & Text & -- & -- & Yes \\
EntityType & Kind of record acted on, e.g. User & Text & -- & -- & Yes \\
EntityId & Id of that record & Text & -- & -- & No \\
Details & Details as JSON & Text & -- & \texttt{\{\}} & Yes \\
CreatedAt & When the action happened & Timestamp & -- & Current time & Yes \\*
Sequence & Position in the hash chain (unique) & Bigint & -- & -- & No \\*
PrevHash & Hash of the previous row & Varchar & 64 & -- & No \\*
Hash & SHA-256 hash of this row & Varchar & 64 & -- & No \\
\end{longtable}
}

There is no separate Admin table: an administrator is a user whose role is Admin or SuperAdmin.

\par\ifdim\dimexpr\pagegoal-\pagetotal\relax<230pt\newpage\fi
\subsection{DeviceBindings Table}
Table~\ref{tab:db_devicebindings} gives the fields of the DeviceBindings table, which links a user to the one PC the account is bound to. Only a hash of the PC's hardware fingerprint is stored.

{\small
\begin{longtable}{>{\raggedright\arraybackslash}p{2.8cm}>{\raggedright\arraybackslash}p{3.7cm}>{\raggedright\arraybackslash}p{1.95cm}>{\centering\arraybackslash}p{1.4cm}>{\raggedright\arraybackslash}p{1.95cm}>{\centering\arraybackslash}p{1.5cm}}
\caption{DeviceBindings Table}\label{tab:db_devicebindings}\\
\hline\hline
\textbf{Field Name} & \textbf{Description} & \textbf{Type} & \textbf{Length} & \textbf{Default} & \textbf{Not Null} \\ \hline
\endfirsthead
\multicolumn{6}{c}{\tablename~\thetable{} -- continued from previous page}\\
\hline\hline
\textbf{Field Name} & \textbf{Description} & \textbf{Type} & \textbf{Length} & \textbf{Default} & \textbf{Not Null} \\ \hline
\endhead
\hline \multicolumn{6}{r}{Continued on next page}\\
\endfoot
\hline\hline
\endlastfoot
BindingId & Unique id of the binding & uuid & -- & Primary Key & Yes \\*
UserId & The bound user (unique; foreign key to Users) & uuid & -- & -- & Yes \\*
HwidHash & SHA-256 hash of the PC's hardware fingerprint & Text & -- & -- & Yes \\
HwidVersion & Fingerprint format: 1 = old, 2 = current & Integer & -- & 1 & Yes \\
MachineName & Windows computer name, e.g. LAB-PC-07 & Text & -- & -- & No \\*
Windows\allowbreak Username & Windows account logged in on that PC & Text & -- & -- & No \\*
RegisteredAt & When the PC was bound & Timestamp & -- & Current time & Yes \\*
LastSeenAt & Last time the bound PC was seen & Timestamp & -- & Current time & Yes \\
\end{longtable}
}

\par\ifdim\dimexpr\pagegoal-\pagetotal\relax<332pt\newpage\fi
\subsection{UserSessions Table}
Table~\ref{tab:db_usersessions} gives the fields of the UserSessions table, which keeps one row for each access token issued at login, so that a token can be revoked before it expires.

{\small
\begin{longtable}{>{\raggedright\arraybackslash}p{2.8cm}>{\raggedright\arraybackslash}p{3.7cm}>{\raggedright\arraybackslash}p{1.95cm}>{\centering\arraybackslash}p{1.4cm}>{\raggedright\arraybackslash}p{1.95cm}>{\centering\arraybackslash}p{1.5cm}}
\caption{UserSessions Table}\label{tab:db_usersessions}\\
\hline\hline
\textbf{Field Name} & \textbf{Description} & \textbf{Type} & \textbf{Length} & \textbf{Default} & \textbf{Not Null} \\ \hline
\endfirsthead
\multicolumn{6}{c}{\tablename~\thetable{} -- continued from previous page}\\
\hline\hline
\textbf{Field Name} & \textbf{Description} & \textbf{Type} & \textbf{Length} & \textbf{Default} & \textbf{Not Null} \\ \hline
\endhead
\hline \multicolumn{6}{r}{Continued on next page}\\
\endfoot
\hline\hline
\endlastfoot
SessionId & Unique id of the row & uuid & -- & Primary Key & Yes \\*
UserId & Owner of the token (foreign key to Users) & uuid & -- & -- & Yes \\*
Jti & JWT ID of the access token (unique) & Text & -- & -- & Yes \\*
ExpiresAt & When the token expires & Timestamp & -- & -- & Yes \\*
IsRevoked & True once the token is revoked & Boolean & -- & false & Yes \\*
CreatedAt & When the token was issued & Timestamp & -- & Current time & Yes \\
\end{longtable}
}

\par\ifdim\dimexpr\pagegoal-\pagetotal\relax<230pt\newpage\fi
\subsection{RefreshTokens Table}
Table~\ref{tab:db_refreshtokens} gives the fields of the RefreshTokens table, which stores the refresh tokens of each user. Only a hash of each token is kept, and every refresh replaces the used token with a new one.

{\small
\begin{longtable}{>{\raggedright\arraybackslash}p{2.8cm}>{\raggedright\arraybackslash}p{3.7cm}>{\raggedright\arraybackslash}p{1.95cm}>{\centering\arraybackslash}p{1.4cm}>{\raggedright\arraybackslash}p{1.95cm}>{\centering\arraybackslash}p{1.5cm}}
\caption{RefreshTokens Table}\label{tab:db_refreshtokens}\\
\hline\hline
\textbf{Field Name} & \textbf{Description} & \textbf{Type} & \textbf{Length} & \textbf{Default} & \textbf{Not Null} \\ \hline
\endfirsthead
\multicolumn{6}{c}{\tablename~\thetable{} -- continued from previous page}\\
\hline\hline
\textbf{Field Name} & \textbf{Description} & \textbf{Type} & \textbf{Length} & \textbf{Default} & \textbf{Not Null} \\ \hline
\endhead
\hline \multicolumn{6}{r}{Continued on next page}\\
\endfoot
\hline\hline
\endlastfoot
TokenId & Unique id of the token & uuid & -- & Primary Key & Yes \\*
UserId & Owner of the token (foreign key to Users) & uuid & -- & -- & Yes \\*
FamilyId & Token family started at login & uuid & -- & -- & Yes \\
TokenHash & SHA-256 hash of the raw token (unique) & Varchar & 64 & -- & Yes \\
Jti & JWT ID of the access token issued with it & Varchar & 64 & -- & Yes \\
Channel & web or student-client & Varchar & 32 & web & Yes \\
MachineName & PC name, sent again as the machine claim on refresh & Varchar & 128 & -- & No \\
CreatedAt & When the token was issued & Timestamp & -- & Current time & Yes \\
ExpiresAt & Fixed expiry; rotation does not extend it & Timestamp & -- & -- & Yes \\
RevokedAt & When the token was revoked (concurrency token) & Timestamp & -- & -- & No \\*
RevokedReason & Why it was revoked, e.g. rotated, reuse or logout & Varchar & 32 & -- & No \\*
Replaced\allowbreak By\allowbreak Token\allowbreak Id & Token that replaced it on rotation (plain id, no foreign key) & uuid & -- & -- & No \\*
CreatedByIp & IP address the token was issued to & Varchar & 64 & -- & No \\
\end{longtable}
}

\par\ifdim\dimexpr\pagegoal-\pagetotal\relax<230pt\newpage\fi
\subsection{Notifications Table}
Table~\ref{tab:db_notifications} gives the fields of the Notifications table, which stores the in-app notifications, such as a new exam or released grades. Each notification has one recipient.

{\small
\begin{longtable}{>{\raggedright\arraybackslash}p{2.8cm}>{\raggedright\arraybackslash}p{3.7cm}>{\raggedright\arraybackslash}p{1.95cm}>{\centering\arraybackslash}p{1.4cm}>{\raggedright\arraybackslash}p{1.95cm}>{\centering\arraybackslash}p{1.5cm}}
\caption{Notifications Table}\label{tab:db_notifications}\\
\hline\hline
\textbf{Field Name} & \textbf{Description} & \textbf{Type} & \textbf{Length} & \textbf{Default} & \textbf{Not Null} \\ \hline
\endfirsthead
\multicolumn{6}{c}{\tablename~\thetable{} -- continued from previous page}\\
\hline\hline
\textbf{Field Name} & \textbf{Description} & \textbf{Type} & \textbf{Length} & \textbf{Default} & \textbf{Not Null} \\ \hline
\endhead
\hline \multicolumn{6}{r}{Continued on next page}\\
\endfoot
\hline\hline
\endlastfoot
NotificationId & Unique id of the notification & uuid & -- & Primary Key & Yes \\*
RecipientId & User who receives it (foreign key to Users) & uuid & -- & -- & Yes \\*
Title & Short title & Text & -- & -- & Yes \\
Body & Message text & Text & -- & -- & Yes \\
Type & Kind, e.g. ExamScheduled or GradeReleased & Text & -- & -- & Yes \\
IsRead & True once the user has read it & Boolean & -- & false & Yes \\*
EmailSent & True if it was also sent by email & Boolean & -- & false & Yes \\*
Related\allowbreak Entity\allowbreak Id & Id of the exam or session it is about & Text & -- & -- & No \\*
CreatedAt & When it was created & Timestamp & -- & Current time & Yes \\
\end{longtable}
}

\par\ifdim\dimexpr\pagegoal-\pagetotal\relax<244pt\newpage\fi
\subsection{Departments Table}
Table~\ref{tab:db_departments} gives the fields of the Departments table, which stores the academic departments. Department names are unique, and a department that still has courses cannot be deleted.

{\small
\begin{longtable}{>{\raggedright\arraybackslash}p{2.8cm}>{\raggedright\arraybackslash}p{3.7cm}>{\raggedright\arraybackslash}p{1.95cm}>{\centering\arraybackslash}p{1.4cm}>{\raggedright\arraybackslash}p{1.95cm}>{\centering\arraybackslash}p{1.5cm}}
\caption{Departments Table}\label{tab:db_departments}\\
\hline\hline
\textbf{Field Name} & \textbf{Description} & \textbf{Type} & \textbf{Length} & \textbf{Default} & \textbf{Not Null} \\ \hline
\endfirsthead
\multicolumn{6}{c}{\tablename~\thetable{} -- continued from previous page}\\
\hline\hline
\textbf{Field Name} & \textbf{Description} & \textbf{Type} & \textbf{Length} & \textbf{Default} & \textbf{Not Null} \\ \hline
\endhead
\hline \multicolumn{6}{r}{Continued on next page}\\
\endfoot
\hline\hline
\endlastfoot
DeptId & Unique id of the department & uuid & -- & Primary Key & Yes \\*
Name & Department name (unique) & Text & -- & -- & Yes \\*
\end{longtable}
}

\par\ifdim\dimexpr\pagegoal-\pagetotal\relax<288pt\newpage\fi
\subsection{Courses Table}
Table~\ref{tab:db_courses} gives the fields of the Courses table, which stores the courses of each department. A course code is unique within its department.

{\small
\begin{longtable}{>{\raggedright\arraybackslash}p{2.8cm}>{\raggedright\arraybackslash}p{3.7cm}>{\raggedright\arraybackslash}p{1.95cm}>{\centering\arraybackslash}p{1.4cm}>{\raggedright\arraybackslash}p{1.95cm}>{\centering\arraybackslash}p{1.5cm}}
\caption{Courses Table}\label{tab:db_courses}\\
\hline\hline
\textbf{Field Name} & \textbf{Description} & \textbf{Type} & \textbf{Length} & \textbf{Default} & \textbf{Not Null} \\ \hline
\endfirsthead
\multicolumn{6}{c}{\tablename~\thetable{} -- continued from previous page}\\
\hline\hline
\textbf{Field Name} & \textbf{Description} & \textbf{Type} & \textbf{Length} & \textbf{Default} & \textbf{Not Null} \\ \hline
\endhead
\hline \multicolumn{6}{r}{Continued on next page}\\
\endfoot
\hline\hline
\endlastfoot
CourseId & Unique id of the course & uuid & -- & Primary Key & Yes \\*
DeptId & Department that offers the course (unique together with Code; foreign key to Departments) & uuid & -- & -- & Yes \\*
Name & Course name & Text & -- & -- & Yes \\*
Code & Course code, e.g. CS301 & Text & -- & -- & Yes \\
\end{longtable}
}

\par\ifdim\dimexpr\pagegoal-\pagetotal\relax<310pt\newpage\fi
\subsection{Sections Table}
Table~\ref{tab:db_sections} gives the fields of the Sections table, which stores a class of a course in one semester, taught by one teacher. Students are enrolled in a section, and exams are created for a section.

{\small
\begin{longtable}{>{\raggedright\arraybackslash}p{2.8cm}>{\raggedright\arraybackslash}p{3.7cm}>{\raggedright\arraybackslash}p{1.95cm}>{\centering\arraybackslash}p{1.4cm}>{\raggedright\arraybackslash}p{1.95cm}>{\centering\arraybackslash}p{1.5cm}}
\caption{Sections Table}\label{tab:db_sections}\\
\hline\hline
\textbf{Field Name} & \textbf{Description} & \textbf{Type} & \textbf{Length} & \textbf{Default} & \textbf{Not Null} \\ \hline
\endfirsthead
\multicolumn{6}{c}{\tablename~\thetable{} -- continued from previous page}\\
\hline\hline
\textbf{Field Name} & \textbf{Description} & \textbf{Type} & \textbf{Length} & \textbf{Default} & \textbf{Not Null} \\ \hline
\endhead
\hline \multicolumn{6}{r}{Continued on next page}\\
\endfoot
\hline\hline
\endlastfoot
SectionId & Unique id of the section & uuid & -- & Primary Key & Yes \\*
CourseId & Course of the section (foreign key to Courses) & uuid & -- & -- & Yes \\*
TeacherId & Teacher of the section (foreign key to Users) & uuid & -- & -- & Yes \\*
Name & Section name & Text & -- & -- & Yes \\*
Semester & Semester of the section & Text & -- & -- & Yes \\
\end{longtable}
}

\par\ifdim\dimexpr\pagegoal-\pagetotal\relax<266pt\newpage\fi
\subsection{SectionEnrollments Table}
Table~\ref{tab:db_sectionenrollments} gives the fields of the SectionEnrollments table, which records which students are enrolled in which section.

{\small
\begin{longtable}{>{\raggedright\arraybackslash}p{2.8cm}>{\raggedright\arraybackslash}p{3.7cm}>{\raggedright\arraybackslash}p{1.95cm}>{\centering\arraybackslash}p{1.4cm}>{\raggedright\arraybackslash}p{1.95cm}>{\centering\arraybackslash}p{1.5cm}}
\caption{SectionEnrollments Table}\label{tab:db_sectionenrollments}\\
\hline\hline
\textbf{Field Name} & \textbf{Description} & \textbf{Type} & \textbf{Length} & \textbf{Default} & \textbf{Not Null} \\ \hline
\endfirsthead
\multicolumn{6}{c}{\tablename~\thetable{} -- continued from previous page}\\
\hline\hline
\textbf{Field Name} & \textbf{Description} & \textbf{Type} & \textbf{Length} & \textbf{Default} & \textbf{Not Null} \\ \hline
\endhead
\hline \multicolumn{6}{r}{Continued on next page}\\
\endfoot
\hline\hline
\endlastfoot
EnrollmentId & Unique id of the enrolment & uuid & -- & Primary Key & Yes \\*
SectionId & The section (foreign key to Sections) & uuid & -- & -- & Yes \\*
UserId & The enrolled student (foreign key to Users) & uuid & -- & -- & Yes \\
\end{longtable}
}

\par\ifdim\dimexpr\pagegoal-\pagetotal\relax<354pt\newpage\fi
\subsection{AttendanceRecords Table}
Table~\ref{tab:db_attendancerecords} gives the fields of the AttendanceRecords table, which stores each student's class attendance in a section, imported from the teacher's CSV file. Exam eligibility is calculated from these rows.

{\small
\begin{longtable}{>{\raggedright\arraybackslash}p{2.8cm}>{\raggedright\arraybackslash}p{3.7cm}>{\raggedright\arraybackslash}p{1.95cm}>{\centering\arraybackslash}p{1.4cm}>{\raggedright\arraybackslash}p{1.95cm}>{\centering\arraybackslash}p{1.5cm}}
\caption{AttendanceRecords Table}\label{tab:db_attendancerecords}\\
\hline\hline
\textbf{Field Name} & \textbf{Description} & \textbf{Type} & \textbf{Length} & \textbf{Default} & \textbf{Not Null} \\ \hline
\endfirsthead
\multicolumn{6}{c}{\tablename~\thetable{} -- continued from previous page}\\
\hline\hline
\textbf{Field Name} & \textbf{Description} & \textbf{Type} & \textbf{Length} & \textbf{Default} & \textbf{Not Null} \\ \hline
\endhead
\hline \multicolumn{6}{r}{Continued on next page}\\
\endfoot
\hline\hline
\endlastfoot
AttendanceId & Unique id of the record & uuid & -- & Primary Key & Yes \\*
SectionId & The section (unique together with UserId; foreign key to Sections) & uuid & -- & -- & Yes \\*
UserId & The student (foreign key to Users) & uuid & -- & -- & Yes \\
TotalClasses & Number of classes held & Integer & -- & -- & Yes \\*
Attended\allowbreak Classes & Number of classes the student attended & Integer & -- & -- & Yes \\*
UpdatedAt & When the row was last imported & Timestamp & -- & Current time & Yes \\*
ImportedBy & User who imported the file (plain id, no foreign key) & uuid & -- & -- & No \\
\end{longtable}
}

\par\ifdim\dimexpr\pagegoal-\pagetotal\relax<288pt\newpage\fi
\subsection{Labs Table}
Table~\ref{tab:db_labs} gives the fields of the Labs table, which stores the computer labs where exams are held. Each lab has many workstations.

{\small
\begin{longtable}{>{\raggedright\arraybackslash}p{2.8cm}>{\raggedright\arraybackslash}p{3.7cm}>{\raggedright\arraybackslash}p{1.95cm}>{\centering\arraybackslash}p{1.4cm}>{\raggedright\arraybackslash}p{1.95cm}>{\centering\arraybackslash}p{1.5cm}}
\caption{Labs Table}\label{tab:db_labs}\\
\hline\hline
\textbf{Field Name} & \textbf{Description} & \textbf{Type} & \textbf{Length} & \textbf{Default} & \textbf{Not Null} \\ \hline
\endfirsthead
\multicolumn{6}{c}{\tablename~\thetable{} -- continued from previous page}\\
\hline\hline
\textbf{Field Name} & \textbf{Description} & \textbf{Type} & \textbf{Length} & \textbf{Default} & \textbf{Not Null} \\ \hline
\endhead
\hline \multicolumn{6}{r}{Continued on next page}\\
\endfoot
\hline\hline
\endlastfoot
LabId & Unique id of the lab & uuid & -- & Primary Key & Yes \\*
Name & Lab name & Text & -- & -- & Yes \\*
Location & Where the lab is & Text & -- & -- & Yes \\*
CreatedAt & When the lab was added & Timestamp & -- & Current time & Yes \\
\end{longtable}
}

\par\ifdim\dimexpr\pagegoal-\pagetotal\relax<310pt\newpage\fi
\subsection{Workstations Table}
Table~\ref{tab:db_workstations} gives the fields of the Workstations table, which stores the lab PCs that are used as exam seats. Each PC belongs to one lab.

{\small
\begin{longtable}{>{\raggedright\arraybackslash}p{2.8cm}>{\raggedright\arraybackslash}p{3.7cm}>{\raggedright\arraybackslash}p{1.95cm}>{\centering\arraybackslash}p{1.4cm}>{\raggedright\arraybackslash}p{1.95cm}>{\centering\arraybackslash}p{1.5cm}}
\caption{Workstations Table}\label{tab:db_workstations}\\
\hline\hline
\textbf{Field Name} & \textbf{Description} & \textbf{Type} & \textbf{Length} & \textbf{Default} & \textbf{Not Null} \\ \hline
\endfirsthead
\multicolumn{6}{c}{\tablename~\thetable{} -- continued from previous page}\\
\hline\hline
\textbf{Field Name} & \textbf{Description} & \textbf{Type} & \textbf{Length} & \textbf{Default} & \textbf{Not Null} \\ \hline
\endhead
\hline \multicolumn{6}{r}{Continued on next page}\\
\endfoot
\hline\hline
\endlastfoot
WorkstationId & Unique id of the PC & uuid & -- & Primary Key & Yes \\*
LabId & Lab the PC is in (foreign key to Labs) & uuid & -- & -- & Yes \\*
MachineNumber & Machine number of the seat & Text & -- & -- & Yes \\*
IpAddress & IP address of the PC & Text & -- & -- & Yes \\*
HostName & Windows computer name, e.g. LAB-PC-07 & Varchar & 128 & -- & No \\
\end{longtable}
}

\par\ifdim\dimexpr\pagegoal-\pagetotal\relax<230pt\newpage\fi
\subsection{Exams Table}
Table~\ref{tab:db_exams} gives the fields of the Exams table, which stores each exam with its schedule and the rules the desktop client enforces. Every exam belongs to one section.

{\small
\begin{longtable}{>{\raggedright\arraybackslash}p{2.8cm}>{\raggedright\arraybackslash}p{3.7cm}>{\raggedright\arraybackslash}p{1.95cm}>{\centering\arraybackslash}p{1.4cm}>{\raggedright\arraybackslash}p{1.95cm}>{\centering\arraybackslash}p{1.5cm}}
\caption{Exams Table}\label{tab:db_exams}\\
\hline\hline
\textbf{Field Name} & \textbf{Description} & \textbf{Type} & \textbf{Length} & \textbf{Default} & \textbf{Not Null} \\ \hline
\endfirsthead
\multicolumn{6}{c}{\tablename~\thetable{} -- continued from previous page}\\
\hline\hline
\textbf{Field Name} & \textbf{Description} & \textbf{Type} & \textbf{Length} & \textbf{Default} & \textbf{Not Null} \\ \hline
\endhead
\hline \multicolumn{6}{r}{Continued on next page}\\
\endfoot
\hline\hline
\endlastfoot
ExamId & Unique id of the exam & uuid & -- & Primary Key & Yes \\*
SectionId & Section the exam is for (foreign key to Sections) & uuid & -- & -- & Yes \\*
Title & Exam title & Text & -- & -- & Yes \\
Instructions & Shown on the student exam dashboard & Text & -- & -- & No \\
StartTime & Scheduled start & Timestamp & -- & -- & Yes \\
Duration\allowbreak Minutes & Length of the exam in minutes & Integer & -- & -- & Yes \\
Early\allowbreak Join\allowbreak Minutes & Minutes before the start that students may begin; 0 = at the start & Integer & -- & -- & Yes \\
AllowedApps & JSON list of allowed programs, e.g. code.exe & Text & -- & \texttt{[]} & Yes \\
Ai\allowbreak Evaluation\allowbreak Enabled & Turns automatic grading on & Boolean & -- & true & Yes \\
Plagiarism\allowbreak Threshold & Similarity percent at which a pair is flagged & Integer & -- & 70 & Yes \\
LockdownMode & Off, Report or Enforce & Varchar & 16 & Report & Yes \\
Block\allowbreak Multi\allowbreak Monitor & Block a second monitor & Boolean & -- & true & Yes \\
Require\allowbreak Readiness & Start needs a passing readiness check at most 30 minutes old & Boolean & -- & -- & Yes \\
RequireWebcam & Exam needs a webcam & Boolean & -- & -- & Yes \\
RequireMic & Exam needs a microphone & Boolean & -- & -- & Yes \\
IsMock & Practice exam: unlimited attempts, left out of results and plagiarism & Boolean & -- & -- & Yes \\
Min\allowbreak Attendance\allowbreak Percent & Attendance needed to be eligible & Double & -- & 75 & Yes \\
EnforceSeat & Student may only sit at the assigned seat & Boolean & -- & true & Yes \\
Screenshot\allowbreak Interval\allowbreak Seconds & Mean time between screenshots & Integer & -- & 60 & Yes \\*
Webcam\allowbreak Interval\allowbreak Seconds & Mean time between webcam snapshots & Integer & -- & 60 & Yes \\*
Status & Scheduled, Active or Ended & Text & -- & Scheduled & Yes \\*
CreatedAt & When the exam was created & Timestamp & -- & Current time & Yes \\
\end{longtable}
}

\par\ifdim\dimexpr\pagegoal-\pagetotal\relax<230pt\newpage\fi
\subsection{ExamAssignments Table}
Table~\ref{tab:db_examassignments} gives the fields of the ExamAssignments table, the roster of an exam. Each row joins a student to an exam and holds the eligibility decision and the optional seat.

{\small
\begin{longtable}{>{\raggedright\arraybackslash}p{2.8cm}>{\raggedright\arraybackslash}p{3.7cm}>{\raggedright\arraybackslash}p{1.95cm}>{\centering\arraybackslash}p{1.4cm}>{\raggedright\arraybackslash}p{1.95cm}>{\centering\arraybackslash}p{1.5cm}}
\caption{ExamAssignments Table}\label{tab:db_examassignments}\\
\hline\hline
\textbf{Field Name} & \textbf{Description} & \textbf{Type} & \textbf{Length} & \textbf{Default} & \textbf{Not Null} \\ \hline
\endfirsthead
\multicolumn{6}{c}{\tablename~\thetable{} -- continued from previous page}\\
\hline\hline
\textbf{Field Name} & \textbf{Description} & \textbf{Type} & \textbf{Length} & \textbf{Default} & \textbf{Not Null} \\ \hline
\endhead
\hline \multicolumn{6}{r}{Continued on next page}\\
\endfoot
\hline\hline
\endlastfoot
AssignmentId & Unique id of the assignment & uuid & -- & Primary Key & Yes \\*
ExamId & The exam (foreign key to Exams) & uuid & -- & -- & Yes \\*
UserId & The student (foreign key to Users) & uuid & -- & -- & Yes \\
WorkstationId & Assigned seat, if any (foreign key to Workstations) & uuid & -- & -- & No \\
IsEligible & Whether the student may sit the exam & Boolean & -- & true & Yes \\*
Eligibility\allowbreak Note & Reason for the eligibility decision & Text & -- & -- & No \\*
Eligibility\allowbreak Overridden & True when a teacher decided by hand & Boolean & -- & -- & Yes \\*
Eligibility\allowbreak Source & Attendance or Manual & Varchar & 16 & -- & No \\
\end{longtable}
}

\par\ifdim\dimexpr\pagegoal-\pagetotal\relax<288pt\newpage\fi
\subsection{ExamProctors Table}
Table~\ref{tab:db_examproctors} gives the fields of the ExamProctors table, which records which staff members supervise which exam. A person can be assigned to an exam only once.

{\small
\begin{longtable}{>{\raggedright\arraybackslash}p{2.8cm}>{\raggedright\arraybackslash}p{3.7cm}>{\raggedright\arraybackslash}p{1.95cm}>{\centering\arraybackslash}p{1.4cm}>{\raggedright\arraybackslash}p{1.95cm}>{\centering\arraybackslash}p{1.5cm}}
\caption{ExamProctors Table}\label{tab:db_examproctors}\\
\hline\hline
\textbf{Field Name} & \textbf{Description} & \textbf{Type} & \textbf{Length} & \textbf{Default} & \textbf{Not Null} \\ \hline
\endfirsthead
\multicolumn{6}{c}{\tablename~\thetable{} -- continued from previous page}\\
\hline\hline
\textbf{Field Name} & \textbf{Description} & \textbf{Type} & \textbf{Length} & \textbf{Default} & \textbf{Not Null} \\ \hline
\endhead
\hline \multicolumn{6}{r}{Continued on next page}\\
\endfoot
\hline\hline
\endlastfoot
ExamId & The exam (foreign key to Exams) & uuid & -- & Primary Key & Yes \\*
UserId & Staff member who supervises it (foreign key to Users) & uuid & -- & Primary Key & Yes \\*
AssignedBy & User who made the assignment (plain id, no foreign key) & uuid & -- & -- & Yes \\*
AssignedAt & When the assignment was made & Timestamp & -- & Current time & Yes \\
\end{longtable}
}

\par\ifdim\dimexpr\pagegoal-\pagetotal\relax<230pt\newpage\fi
\subsection{Questions Table}
Table~\ref{tab:db_questions} gives the fields of the Questions table, which stores the coding, theory and multiple-choice questions of each exam.

{\small
\begin{longtable}{>{\raggedright\arraybackslash}p{2.8cm}>{\raggedright\arraybackslash}p{3.7cm}>{\raggedright\arraybackslash}p{1.95cm}>{\centering\arraybackslash}p{1.4cm}>{\raggedright\arraybackslash}p{1.95cm}>{\centering\arraybackslash}p{1.5cm}}
\caption{Questions Table}\label{tab:db_questions}\\
\hline\hline
\textbf{Field Name} & \textbf{Description} & \textbf{Type} & \textbf{Length} & \textbf{Default} & \textbf{Not Null} \\ \hline
\endfirsthead
\multicolumn{6}{c}{\tablename~\thetable{} -- continued from previous page}\\
\hline\hline
\textbf{Field Name} & \textbf{Description} & \textbf{Type} & \textbf{Length} & \textbf{Default} & \textbf{Not Null} \\ \hline
\endhead
\hline \multicolumn{6}{r}{Continued on next page}\\
\endfoot
\hline\hline
\endlastfoot
QuestionId & Unique id of the question & uuid & -- & Primary Key & Yes \\*
ExamId & Exam the question is in (foreign key to Exams) & uuid & -- & -- & Yes \\*
Type & Coding, Theory or Mcq & Text & -- & -- & Yes \\
BodyText & Question text & Text & -- & -- & Yes \\
Marks & Marks for the question & Integer & -- & -- & Yes \\
ExpectedOutput & Expected program output & Text & -- & -- & No \\
OrderIndex & Position in the exam & Integer & -- & -- & Yes \\
OptionsJson & MCQ only: options as a JSON list & Text & -- & -- & No \\
Correct\allowbreak Option\allowbreak Index & MCQ only: index of the correct option; never sent to students & Integer & -- & -- & No \\*
Reference\allowbreak Answer & Model answer or sample solution; never sent to students & Text & -- & -- & No \\*
Language & Coding: cpp or c; null = detect from the code & Varchar & 16 & -- & No \\*
RubricKeywords & Theory: key terms the grader looks for, comma-separated & Varchar & 1000 & -- & No \\
\end{longtable}
}

\par\ifdim\dimexpr\pagegoal-\pagetotal\relax<230pt\newpage\fi
\subsection{TestCases Table}
Table~\ref{tab:db_testcases} gives the fields of the TestCases table, which stores the test cases of a coding question. The weight of a case decides its share of the marks, and hidden cases are not shown to students.

{\small
\begin{longtable}{>{\raggedright\arraybackslash}p{2.8cm}>{\raggedright\arraybackslash}p{3.7cm}>{\raggedright\arraybackslash}p{1.95cm}>{\centering\arraybackslash}p{1.4cm}>{\raggedright\arraybackslash}p{1.95cm}>{\centering\arraybackslash}p{1.5cm}}
\caption{TestCases Table}\label{tab:db_testcases}\\
\hline\hline
\textbf{Field Name} & \textbf{Description} & \textbf{Type} & \textbf{Length} & \textbf{Default} & \textbf{Not Null} \\ \hline
\endfirsthead
\multicolumn{6}{c}{\tablename~\thetable{} -- continued from previous page}\\
\hline\hline
\textbf{Field Name} & \textbf{Description} & \textbf{Type} & \textbf{Length} & \textbf{Default} & \textbf{Not Null} \\ \hline
\endhead
\hline \multicolumn{6}{r}{Continued on next page}\\
\endfoot
\hline\hline
\endlastfoot
TestCaseId & Unique id of the test case & uuid & -- & Primary Key & Yes \\*
QuestionId & Question the case belongs to (foreign key to Questions) & uuid & -- & -- & Yes \\*
Input & Input given to the program & Text & -- & -- & Yes \\
ExpectedOutput & Output the program must print & Text & -- & -- & Yes \\
IsHidden & Hidden from students & Boolean & -- & false & Yes \\
Weight & Share of the question's marks & Double & -- & 1 & Yes \\
TimeLimitMs & Time limit for one run; null = the 2-second default & Integer & -- & -- & No \\
IsEdgeCase & Marks the case as an edge case & Boolean & -- & -- & Yes \\
Category & Free label, e.g. basic & Varchar & 32 & -- & No \\*
IgnoreCase & Ignore letter case when comparing output & Boolean & -- & -- & Yes \\*
FloatTolerance & Allowed difference for numbers; null = exact text & Double & -- & -- & No \\*
OrderIndex & Order inside the question & Integer & -- & -- & Yes \\
\end{longtable}
}

\par\ifdim\dimexpr\pagegoal-\pagetotal\relax<354pt\newpage\fi
\subsection{ReadinessChecks Table}
Table~\ref{tab:db_readinesschecks} gives the fields of the ReadinessChecks table, which stores each readiness check the desktop client reports before an exam. When the exam requires the check, the latest one decides whether the student may start.

{\small
\begin{longtable}{>{\raggedright\arraybackslash}p{2.8cm}>{\raggedright\arraybackslash}p{3.7cm}>{\raggedright\arraybackslash}p{1.95cm}>{\centering\arraybackslash}p{1.4cm}>{\raggedright\arraybackslash}p{1.95cm}>{\centering\arraybackslash}p{1.5cm}}
\caption{ReadinessChecks Table}\label{tab:db_readinesschecks}\\
\hline\hline
\textbf{Field Name} & \textbf{Description} & \textbf{Type} & \textbf{Length} & \textbf{Default} & \textbf{Not Null} \\ \hline
\endfirsthead
\multicolumn{6}{c}{\tablename~\thetable{} -- continued from previous page}\\
\hline\hline
\textbf{Field Name} & \textbf{Description} & \textbf{Type} & \textbf{Length} & \textbf{Default} & \textbf{Not Null} \\ \hline
\endhead
\hline \multicolumn{6}{r}{Continued on next page}\\
\endfoot
\hline\hline
\endlastfoot
Id & Unique id of the check & uuid & -- & Primary Key & Yes \\*
ExamId & The exam (foreign key to Exams) & uuid & -- & -- & Yes \\*
UserId & The student (foreign key to Users) & uuid & -- & -- & Yes \\
PassedAll & True if every check passed & Boolean & -- & -- & Yes \\*
ResultsJson & Name, result, detail and time of each check, as JSON & Text & -- & \texttt{[]} & Yes \\*
MachineName & PC the check ran on & Varchar & 128 & -- & No \\*
CheckedAt & When the check ran & Timestamp & -- & Current time & Yes \\
\end{longtable}
}

\par\ifdim\dimexpr\pagegoal-\pagetotal\relax<230pt\newpage\fi
\subsection{ExamSessions Table}
Table~\ref{tab:db_examsessions} gives the fields of the ExamSessions table, which stores one student's attempt at one exam. The answers, monitoring events, screenshots and incidents of the attempt point to it.

{\small
\begin{longtable}{>{\raggedright\arraybackslash}p{2.8cm}>{\raggedright\arraybackslash}p{3.7cm}>{\raggedright\arraybackslash}p{1.95cm}>{\centering\arraybackslash}p{1.4cm}>{\raggedright\arraybackslash}p{1.95cm}>{\centering\arraybackslash}p{1.5cm}}
\caption{ExamSessions Table}\label{tab:db_examsessions}\\
\hline\hline
\textbf{Field Name} & \textbf{Description} & \textbf{Type} & \textbf{Length} & \textbf{Default} & \textbf{Not Null} \\ \hline
\endfirsthead
\multicolumn{6}{c}{\tablename~\thetable{} -- continued from previous page}\\
\hline\hline
\textbf{Field Name} & \textbf{Description} & \textbf{Type} & \textbf{Length} & \textbf{Default} & \textbf{Not Null} \\ \hline
\endhead
\hline \multicolumn{6}{r}{Continued on next page}\\
\endfoot
\hline\hline
\endlastfoot
SessionId & Unique id of the attempt & uuid & -- & Primary Key & Yes \\*
ExamId & The exam (foreign key to Exams) & uuid & -- & -- & Yes \\*
UserId & The student (foreign key to Users) & uuid & -- & -- & Yes \\
StartedAt & When the attempt started & Timestamp & -- & Current time & Yes \\
SubmittedAt & When the attempt was closed & Timestamp & -- & -- & No \\
Status & InProgress, Submitted, ForceSubmitted or TimedOut (concurrency token) & Text & -- & InProgress & Yes \\
ExtraMinutes & Extra time given by a proctor & Integer & -- & -- & Yes \\
ReceiptSha256 & SHA-256 receipt over the stored answers & Varchar & 64 & -- & No \\
Receipt\allowbreak Answer\allowbreak Count & Number of answers the receipt covers & Integer & -- & -- & No \\
Last\allowbreak Heartbeat\allowbreak At & Time of the latest heartbeat & Timestamp & -- & -- & No \\*
IsDisconnected & True while the client is silent & Boolean & -- & -- & Yes \\*
DisconnectedAt & When the client was marked disconnected & Timestamp & -- & -- & No \\*
UnlockedUntil & End of a proctor's lockdown unlock & Timestamp & -- & -- & No \\
\end{longtable}
}

\par\ifdim\dimexpr\pagegoal-\pagetotal\relax<230pt\newpage\fi
\subsection{Answers Table}
Table~\ref{tab:db_answers} gives the fields of the Answers table, which stores the latest answer of a session to each question. A revision number stops an older save that arrives late from overwriting a newer answer.

{\small
\begin{longtable}{>{\raggedright\arraybackslash}p{2.8cm}>{\raggedright\arraybackslash}p{3.7cm}>{\raggedright\arraybackslash}p{1.95cm}>{\centering\arraybackslash}p{1.4cm}>{\raggedright\arraybackslash}p{1.95cm}>{\centering\arraybackslash}p{1.5cm}}
\caption{Answers Table}\label{tab:db_answers}\\
\hline\hline
\textbf{Field Name} & \textbf{Description} & \textbf{Type} & \textbf{Length} & \textbf{Default} & \textbf{Not Null} \\ \hline
\endfirsthead
\multicolumn{6}{c}{\tablename~\thetable{} -- continued from previous page}\\
\hline\hline
\textbf{Field Name} & \textbf{Description} & \textbf{Type} & \textbf{Length} & \textbf{Default} & \textbf{Not Null} \\ \hline
\endhead
\hline \multicolumn{6}{r}{Continued on next page}\\
\endfoot
\hline\hline
\endlastfoot
AnswerId & Unique id of the answer & uuid & -- & Primary Key & Yes \\*
SessionId & The attempt (unique together with QuestionId; foreign key to ExamSessions) & uuid & -- & -- & Yes \\*
QuestionId & The question answered (foreign key to Questions) & uuid & -- & -- & Yes \\
AnswerText & The student's answer & Text & -- & -- & Yes \\
SubmittedAt & When the answer was submitted & Timestamp & -- & -- & No \\*
LastSavedAt & Last time the server saved it & Timestamp & -- & Current time & Yes \\*
Revision & Edit counter from the client; a lower revision is refused (concurrency token) & Integer & -- & -- & Yes \\*
ClientSavedAt & Save time reported by the client, for diagnostics & Timestamp & -- & -- & No \\
\end{longtable}
}

\par\ifdim\dimexpr\pagegoal-\pagetotal\relax<354pt\newpage\fi
\subsection{MonitoringEvents Table}
Table~\ref{tab:db_monitoringevents} gives the fields of the MonitoringEvents table, which stores the events of an exam session, such as heartbeats and violations. An event sent twice in an offline batch is stored only once.

{\small
\begin{longtable}{>{\raggedright\arraybackslash}p{2.8cm}>{\raggedright\arraybackslash}p{3.7cm}>{\raggedright\arraybackslash}p{1.95cm}>{\centering\arraybackslash}p{1.4cm}>{\raggedright\arraybackslash}p{1.95cm}>{\centering\arraybackslash}p{1.5cm}}
\caption{MonitoringEvents Table}\label{tab:db_monitoringevents}\\
\hline\hline
\textbf{Field Name} & \textbf{Description} & \textbf{Type} & \textbf{Length} & \textbf{Default} & \textbf{Not Null} \\ \hline
\endfirsthead
\multicolumn{6}{c}{\tablename~\thetable{} -- continued from previous page}\\
\hline\hline
\textbf{Field Name} & \textbf{Description} & \textbf{Type} & \textbf{Length} & \textbf{Default} & \textbf{Not Null} \\ \hline
\endhead
\hline \multicolumn{6}{r}{Continued on next page}\\
\endfoot
\hline\hline
\endlastfoot
EventId & Unique id of the event & uuid & -- & Primary Key & Yes \\*
ExamSessionId & The attempt (unique together with ClientEventId when set; foreign key to ExamSessions) & uuid & -- & -- & Yes \\*
EventType & Kind of event, e.g. FocusLost or PhoneDetected & Text & -- & -- & Yes \\
Payload & Event details as JSON & Text & -- & \texttt{\{\}} & Yes \\*
RecordedAt & Time on the student's PC & Timestamp & -- & Current time & Yes \\*
ClientEventId & Id the client gave the event & uuid & -- & -- & No \\*
ReceivedAt & When the server received it & Timestamp & -- & -- & No \\
\end{longtable}
}

\par\ifdim\dimexpr\pagegoal-\pagetotal\relax<230pt\newpage\fi
\subsection{Screenshots Table}
Table~\ref{tab:db_screenshots} gives the fields of the Screenshots table, which stores the details of each screen capture or webcam snapshot. The image is kept on the server's disk with a hash, so it can be checked later.

{\small
\begin{longtable}{>{\raggedright\arraybackslash}p{2.8cm}>{\raggedright\arraybackslash}p{3.7cm}>{\raggedright\arraybackslash}p{1.95cm}>{\centering\arraybackslash}p{1.4cm}>{\raggedright\arraybackslash}p{1.95cm}>{\centering\arraybackslash}p{1.5cm}}
\caption{Screenshots Table}\label{tab:db_screenshots}\\
\hline\hline
\textbf{Field Name} & \textbf{Description} & \textbf{Type} & \textbf{Length} & \textbf{Default} & \textbf{Not Null} \\ \hline
\endfirsthead
\multicolumn{6}{c}{\tablename~\thetable{} -- continued from previous page}\\
\hline\hline
\textbf{Field Name} & \textbf{Description} & \textbf{Type} & \textbf{Length} & \textbf{Default} & \textbf{Not Null} \\ \hline
\endhead
\hline \multicolumn{6}{r}{Continued on next page}\\
\endfoot
\hline\hline
\endlastfoot
ScreenshotId & Unique id of the capture & uuid & -- & Primary Key & Yes \\*
SessionId & The attempt (foreign key to ExamSessions) & uuid & -- & -- & Yes \\*
ExamId & Exam of the attempt (plain id, no foreign key) & uuid & -- & -- & Yes \\
Kind & Screen or Webcam & Varchar & 16 & Screen & Yes \\
Trigger & Interval, Violation or Manual & Varchar & 16 & Interval & Yes \\
TriggerReason & Violation that caused it, e.g. UnauthorizedProcess & Varchar & 64 & -- & No \\
CapturedAt & Capture time on the client & Timestamp & -- & -- & Yes \\
UploadedAt & When the server received it & Timestamp & -- & Current time & Yes \\
StoragePath & Path of the image file on the server & Varchar & 512 & -- & Yes \\*
ContentType & Image type & Varchar & 32 & image/jpeg & Yes \\*
SizeBytes & File size in bytes & Integer & -- & -- & Yes \\*
Sha256 & SHA-256 hash of the stored file & Varchar & 64 & -- & Yes \\
\end{longtable}
}

\par\ifdim\dimexpr\pagegoal-\pagetotal\relax<230pt\newpage\fi
\subsection{Incidents Table}
Table~\ref{tab:db_incidents} gives the fields of the Incidents table, which stores what an invigilator saw and recorded during a session. An incident can be escalated to the teacher and resolved with a note.

{\small
\begin{longtable}{>{\raggedright\arraybackslash}p{2.8cm}>{\raggedright\arraybackslash}p{3.7cm}>{\raggedright\arraybackslash}p{1.95cm}>{\centering\arraybackslash}p{1.4cm}>{\raggedright\arraybackslash}p{1.95cm}>{\centering\arraybackslash}p{1.5cm}}
\caption{Incidents Table}\label{tab:db_incidents}\\
\hline\hline
\textbf{Field Name} & \textbf{Description} & \textbf{Type} & \textbf{Length} & \textbf{Default} & \textbf{Not Null} \\ \hline
\endfirsthead
\multicolumn{6}{c}{\tablename~\thetable{} -- continued from previous page}\\
\hline\hline
\textbf{Field Name} & \textbf{Description} & \textbf{Type} & \textbf{Length} & \textbf{Default} & \textbf{Not Null} \\ \hline
\endhead
\hline \multicolumn{6}{r}{Continued on next page}\\
\endfoot
\hline\hline
\endlastfoot
IncidentId & Unique id of the incident & uuid & -- & Primary Key & Yes \\*
SessionId & The attempt (foreign key to ExamSessions) & uuid & -- & -- & Yes \\*
ExamId & Exam of the attempt (plain id, no foreign key) & uuid & -- & -- & Yes \\
CreatedBy & Staff member who recorded it (plain id, no foreign key) & uuid & -- & -- & Yes \\
Severity & Low, Medium, High or Critical & Varchar & 16 & Medium & Yes \\
Status & Open, Escalated or Resolved & Varchar & 16 & Open & Yes \\
Notes & What the proctor saw & Varchar & 2000 & -- & Yes \\
StatusNote & Note given with the latest status change & Varchar & 2000 & -- & No \\
CreatedAt & When it was recorded & Timestamp & -- & Current time & Yes \\*
UpdatedAt & When it last changed & Timestamp & -- & Current time & Yes \\*
ResolvedAt & When it was resolved & Timestamp & -- & -- & No \\*
ResolvedBy & Who resolved it (plain id, no foreign key) & uuid & -- & -- & No \\
\end{longtable}
}

\par\ifdim\dimexpr\pagegoal-\pagetotal\relax<230pt\newpage\fi
\subsection{AiGradingResults Table}
Table~\ref{tab:db_aigradingresults} gives the fields of the AiGradingResults table, which stores the suggested mark for one answer. A teacher's override always takes precedence over it.

{\small
\begin{longtable}{>{\raggedright\arraybackslash}p{2.8cm}>{\raggedright\arraybackslash}p{3.7cm}>{\raggedright\arraybackslash}p{1.95cm}>{\centering\arraybackslash}p{1.4cm}>{\raggedright\arraybackslash}p{1.95cm}>{\centering\arraybackslash}p{1.5cm}}
\caption{AiGradingResults Table}\label{tab:db_aigradingresults}\\
\hline\hline
\textbf{Field Name} & \textbf{Description} & \textbf{Type} & \textbf{Length} & \textbf{Default} & \textbf{Not Null} \\ \hline
\endfirsthead
\multicolumn{6}{c}{\tablename~\thetable{} -- continued from previous page}\\
\hline\hline
\textbf{Field Name} & \textbf{Description} & \textbf{Type} & \textbf{Length} & \textbf{Default} & \textbf{Not Null} \\ \hline
\endhead
\hline \multicolumn{6}{r}{Continued on next page}\\
\endfoot
\hline\hline
\endlastfoot
ResultId & Unique id of the result & uuid & -- & Primary Key & Yes \\*
AnswerId & The graded answer (unique; foreign key to Answers) & uuid & -- & -- & Yes \\*
SuggestedMarks & Suggested mark & Double & -- & -- & Yes \\
Justification & Reason for the mark, shown to the teacher & Text & -- & -- & Yes \\
Confidence & High, Medium or Low & Text & -- & Medium & Yes \\
GradedAt & When it was graded & Timestamp & -- & Current time & Yes \\
Status & Pending, Graded or Failed & Varchar & 16 & Graded & Yes \\
Engine & What produced it, e.g. MCQ or TF-IDF & Varchar & 64 & -- & No \\
CompileOutput & Compiler messages (shortened), teacher only & Varchar & 4104 & -- & No \\
DurationMs & Time the sandbox run took & Integer & -- & -- & No \\*
AnalysisJson & Code analysis as JSON, teacher only & Text & -- & -- & No \\*
LlmSummary & Optional LLM summary; off by default & Text & -- & -- & No \\*
QueueVersion & Raised on each regrade so an older run cannot overwrite it (concurrency token) & Integer & -- & -- & Yes \\
\end{longtable}
}

\par\ifdim\dimexpr\pagegoal-\pagetotal\relax<230pt\newpage\fi
\subsection{TestCaseResults Table}
Table~\ref{tab:db_testcaseresults} gives the fields of the TestCaseResults table, which stores the outcome of one test case for one graded coding answer.

{\small
\begin{longtable}{>{\raggedright\arraybackslash}p{2.8cm}>{\raggedright\arraybackslash}p{3.7cm}>{\raggedright\arraybackslash}p{1.95cm}>{\centering\arraybackslash}p{1.4cm}>{\raggedright\arraybackslash}p{1.95cm}>{\centering\arraybackslash}p{1.5cm}}
\caption{TestCaseResults Table}\label{tab:db_testcaseresults}\\
\hline\hline
\textbf{Field Name} & \textbf{Description} & \textbf{Type} & \textbf{Length} & \textbf{Default} & \textbf{Not Null} \\ \hline
\endfirsthead
\multicolumn{6}{c}{\tablename~\thetable{} -- continued from previous page}\\
\hline\hline
\textbf{Field Name} & \textbf{Description} & \textbf{Type} & \textbf{Length} & \textbf{Default} & \textbf{Not Null} \\ \hline
\endhead
\hline \multicolumn{6}{r}{Continued on next page}\\
\endfoot
\hline\hline
\endlastfoot
Test\allowbreak Case\allowbreak Result\allowbreak Id & Unique id of the outcome & uuid & -- & Primary Key & Yes \\*
ResultId & The grading result (foreign key to AiGradingResults) & uuid & -- & -- & Yes \\*
TestCaseId & Source test case; null for the compile-only case (plain id, no foreign key) & uuid & -- & -- & No \\
OrderIndex & Order of the case & Integer & -- & -- & Yes \\
Name & Name of the case & Varchar & 64 & -- & Yes \\
Input & Input given to the program & Text & -- & -- & Yes \\
ExpectedOutput & Output the program had to print & Text & -- & -- & Yes \\
IsHidden & Hidden from students & Boolean & -- & -- & Yes \\
IsEdgeCase & Edge case & Boolean & -- & -- & Yes \\
Category & Label of the case & Varchar & 32 & -- & No \\
Weight & Weight of the case & Double & -- & 1 & Yes \\
Status & Pending, Passed, Failed, TimedOut, MemoryExceeded, RuntimeError or NotRun & Varchar & 16 & Pending & Yes \\
Passed & True if the case passed & Boolean & -- & -- & Yes \\
ActualOutput & Program output (shortened) & Varchar & 4104 & -- & No \\*
ExitCode & Exit code of the program & Integer & -- & -- & No \\*
DurationMs & Run time of the case & Integer & -- & -- & No \\*
CompareMode & How the output matched when it passed & Varchar & 32 & -- & No \\
\end{longtable}
}

\par\ifdim\dimexpr\pagegoal-\pagetotal\relax<332pt\newpage\fi
\subsection{TeacherGradeOverrides Table}
Table~\ref{tab:db_teachergradeoverrides} gives the fields of the TeacherGradeOverrides table, which stores the final mark a teacher sets for one answer. It is used instead of the suggested mark.

{\small
\begin{longtable}{>{\raggedright\arraybackslash}p{2.8cm}>{\raggedright\arraybackslash}p{3.7cm}>{\raggedright\arraybackslash}p{1.95cm}>{\centering\arraybackslash}p{1.4cm}>{\raggedright\arraybackslash}p{1.95cm}>{\centering\arraybackslash}p{1.5cm}}
\caption{TeacherGradeOverrides Table}\label{tab:db_teachergradeoverrides}\\
\hline\hline
\textbf{Field Name} & \textbf{Description} & \textbf{Type} & \textbf{Length} & \textbf{Default} & \textbf{Not Null} \\ \hline
\endfirsthead
\multicolumn{6}{c}{\tablename~\thetable{} -- continued from previous page}\\
\hline\hline
\textbf{Field Name} & \textbf{Description} & \textbf{Type} & \textbf{Length} & \textbf{Default} & \textbf{Not Null} \\ \hline
\endhead
\hline \multicolumn{6}{r}{Continued on next page}\\
\endfoot
\hline\hline
\endlastfoot
OverrideId & Unique id of the override & uuid & -- & Primary Key & Yes \\*
AnswerId & The answer (unique; foreign key to Answers) & uuid & -- & -- & Yes \\*
TeacherId & Teacher who set the mark (foreign key to Users) & uuid & -- & -- & Yes \\*
FinalMarks & Final mark & Double & -- & -- & Yes \\*
Note & Teacher's note & Text & -- & -- & No \\*
OverriddenAt & When the mark was set & Timestamp & -- & Current time & Yes \\
\end{longtable}
}

\par\ifdim\dimexpr\pagegoal-\pagetotal\relax<230pt\newpage\fi
\subsection{PlagiarismResults Table}
Table~\ref{tab:db_plagiarismresults} gives the fields of the PlagiarismResults table, which stores pairs of students with similar answers to the same question. Only pairs at or above the exam's threshold are stored.

{\small
\begin{longtable}{>{\raggedright\arraybackslash}p{2.8cm}>{\raggedright\arraybackslash}p{3.7cm}>{\raggedright\arraybackslash}p{1.95cm}>{\centering\arraybackslash}p{1.4cm}>{\raggedright\arraybackslash}p{1.95cm}>{\centering\arraybackslash}p{1.5cm}}
\caption{PlagiarismResults Table}\label{tab:db_plagiarismresults}\\
\hline\hline
\textbf{Field Name} & \textbf{Description} & \textbf{Type} & \textbf{Length} & \textbf{Default} & \textbf{Not Null} \\ \hline
\endfirsthead
\multicolumn{6}{c}{\tablename~\thetable{} -- continued from previous page}\\
\hline\hline
\textbf{Field Name} & \textbf{Description} & \textbf{Type} & \textbf{Length} & \textbf{Default} & \textbf{Not Null} \\ \hline
\endhead
\hline \multicolumn{6}{r}{Continued on next page}\\
\endfoot
\hline\hline
\endlastfoot
PlagId & Unique id of the result & uuid & -- & Primary Key & Yes \\*
ExamId & The exam (foreign key to Exams) & uuid & -- & -- & Yes \\*
QuestionId & The question compared (foreign key to Questions) & uuid & -- & -- & Yes \\
UserIdA & First student of the pair (foreign key to Users) & uuid & -- & -- & Yes \\
UserIdB & Second student of the pair (foreign key to Users) & uuid & -- & -- & Yes \\*
Similarity\allowbreak Score & Similarity in percent & Double & -- & -- & Yes \\*
Matching\allowbreak Segments & Matching parts as JSON & Text & -- & \texttt{[]} & Yes \\*
DetectedAt & When the pair was found & Timestamp & -- & Current time & Yes \\
\end{longtable}
}

\par\ifdim\dimexpr\pagegoal-\pagetotal\relax<332pt\newpage\fi
\subsection{ExamReports Table}
Table~\ref{tab:db_examreports} gives the fields of the ExamReports table, which records each report generated for an exam. The file is saved on the server's disk, and the row keeps its path.

{\small
\begin{longtable}{>{\raggedright\arraybackslash}p{2.8cm}>{\raggedright\arraybackslash}p{3.7cm}>{\raggedright\arraybackslash}p{1.95cm}>{\centering\arraybackslash}p{1.4cm}>{\raggedright\arraybackslash}p{1.95cm}>{\centering\arraybackslash}p{1.5cm}}
\caption{ExamReports Table}\label{tab:db_examreports}\\
\hline\hline
\textbf{Field Name} & \textbf{Description} & \textbf{Type} & \textbf{Length} & \textbf{Default} & \textbf{Not Null} \\ \hline
\endfirsthead
\multicolumn{6}{c}{\tablename~\thetable{} -- continued from previous page}\\
\hline\hline
\textbf{Field Name} & \textbf{Description} & \textbf{Type} & \textbf{Length} & \textbf{Default} & \textbf{Not Null} \\ \hline
\endhead
\hline \multicolumn{6}{r}{Continued on next page}\\
\endfoot
\hline\hline
\endlastfoot
ReportId & Unique id of the report & uuid & -- & Primary Key & Yes \\*
ExamId & The exam (foreign key to Exams) & uuid & -- & -- & Yes \\*
GeneratedBy & User who asked for it (foreign key to Users) & uuid & -- & -- & Yes \\*
ReportType & Kind, e.g. Scores, Attendance or EvidenceZip & Text & -- & -- & Yes \\*
FilePath & Path of the file on the server & Text & -- & -- & Yes \\*
GeneratedAt & When it was made & Timestamp & -- & Current time & Yes \\
\end{longtable}
}

\emergencystretch=\BTsavedES\relax
